\documentclass{article}
\usepackage[T2A]{fontenc}
\usepackage[utf8]{inputenc}   
\usepackage[english, russian]{babel}

% Set page size and margins
% Replace `letterpaper' with`a4paper' for UK/EU standard size
\usepackage[a4paper,top=2cm,bottom=2cm,left=2cm,right=2cm,marginparwidth=1.75cm]{geometry}

\usepackage{amsmath}
\usepackage{graphicx}
\usepackage[colorlinks=true, allcolors=blue]{hyperref}
\usepackage{amsfonts}
\usepackage{amssymb}
% \usepackage[left=1cm,right=1cm,top=1cm,bottom=1cm]{geometry}
\usepackage{hyperref}
\usepackage{seqsplit}
\usepackage[dvipsnames]{xcolor}
\usepackage{enumitem}
\usepackage{algorithm}
\usepackage{algpseudocode}
\usepackage{algorithmicx}
\usepackage{mathalfa}
\usepackage{mathrsfs}
\usepackage{dsfont}
\usepackage{caption,subcaption}
\usepackage{wrapfig}
\usepackage[stable]{footmisc}
\usepackage{indentfirst}
\usepackage{rotating}
\usepackage{pdflscape}
\usepackage{listings}
\usepackage{amsmath}
\usepackage{graphicx}
\usepackage[colorlinks=true, allcolors=blue]{hyperref}
\usepackage{amsfonts}
\usepackage{amssymb}
% \usepackage[left=1cm,right=1cm,top=1cm,bottom=1cm]{geometry}
\usepackage{hyperref}
\usepackage{seqsplit}
\usepackage[dvipsnames]{xcolor}
\usepackage{enumitem}
\usepackage{algorithm}
\usepackage{algpseudocode}
\usepackage{algorithmicx}
\usepackage{mathalfa}
\usepackage{mathrsfs}
\usepackage{dsfont}
\usepackage{caption,subcaption}
\usepackage{wrapfig}
\usepackage[stable]{footmisc}
\usepackage{indentfirst}
\usepackage{rotating}
\usepackage{pdflscape}
\usepackage{graphicx}

\usepackage{minted}

\usepackage{MnSymbol,wasysym}
\usepackage{tikz}

\usepackage{minted}

\usepackage{MnSymbol,wasysym}

\title{Домашнее задание 12}
\author{Лиза Фоменко}
\date{}

\begin{document}

\maketitle


\section*{\textbf{Задание 1}} \textbf{\textit{Примените алгоритм Беллмана-Форда к графу для поиска кратчайших путей от вершины $e$ до всех остальных}}

\begin{figure}[H]
     \centering
     \begin{subfigure}[b]{0.95\textwidth}
         \centering
         \includegraphics[width=\textwidth]{algos/bf.png}
         \caption{Bellman-Ford algorithm}
     \end{subfigure}
     \hfill
\end{figure}

\bigskip

\section*{\textbf{Задание 2}} \textbf{\textit{Пусть сильно связный граф, в котором кратчайшие расстояния между всеми парами вершин имеют реберную длину не более $k$, называется \textit{k-плотненьким}.
Предложите эффективный алгоритм поиска кратчайших путей в \textit{k-плотненьких} графах. Оцените его асимптотику.}}


\bigskip


\section*{\textbf{Задание 3}} \textbf{\textit{Постройте алгоритм, определяющий, содержит ли данный неориентированный граф простой цикл длины 4. Решите в двух случаях, когда имеется в виду реберная длина и вес пути.
}}

\textbf{а) реберная длина}

$v_1 - v_2 - v_3 - v_4 -v_1$ - цикл, который мы ищем. Заметим, что $v_1$ связана с $v_2, v_4$ и $v_3$ связана с $v_2, v_4$. Значит, если мы найдем такие две вершины, у которых есть 2 общих соседа, то в графе есть путь длины 4. На самом деле может быть и больше 2 общих соседей (например у нас квадратный цикл и еще общая вершина $v_1, v_3$, она все равно не нарушает цикл. Значит, нам нужно найти такие вершины, у которых больше 1 общего соседа.

Для этого нам нужно уметь сравнивать списки смердженных вершин для любых двух пар. Например, пусть у нас даны списки смежности 
$v = [x_1, ... x_i]$
$u = [y_1, ... y_j]$

Тогда мы можем сделать так:

\begin{verbatim}
counts = {}

for x in v_merged:
    counts[x] = 1

k = 0                       # счетчик общих вершин
for y in u:
    if y in counts.keys():
        k += 1

if k>1:
    print('got cycle of len 4')
\end{verbatim}

Если у нас общих вершин у $v$ и $u$ больше 1, значит есть цикл длины 4. 
Цикл работает за $O(|x| + |y|)$, где $|x|$ = количество смердженных вершин с $v$, $|y|$ - с $u$. 

Но нам придется проверить попарно все вершины $u, v$ - то есть $O(|V|^2)$ пар. Для фиксированной вершины мы переберем все ребра графа (когда фиксируем $v$ и перебираем остальные $|V| - 1$ вершину, мы делаем проход по всем ребрам графа). Тогда общая ассимптотика - $O(|V| * |V-1|*|E|) = O(|V|^2|E|)$



\textbf{б) вес пути}

НАсколько я помню, поиск цикла заданного веса в орграфе является NP полной задачей, то есть не решается за полиномиальное время. Скорее всего это выводится из гамильтонова цикла. Соответственно наше время будет зависеть от $2^{|V|}$. 
\bigskip

\section*{\textbf{Задание 4}} \textbf{\textit{Дан ориентированный граф на $n$ вершинах $(V = {1, \dots,n})$, который получен из графа-пути (рёбра которого ведут из вершины $i$ в $i + 1$) добавлением ещё каких-то $m$ данных ребер. Найдите количество сильно связных компонент в этом графе за $O(m \log m)$}}


\bigskip

\section*{\textbf{Задание 5}} \textbf{\textit{В орграфе есть ребра отрицательного веса, но нет циклов с отрицательным весом. Предложите алгоритм, который находит для данной вершины вершину, \textbf{от которой} она удалена на максимальное расстояние. Докажите его корректность и оцените асимптотику.}}

Я так понимаю, нам нужно найти максимальное из минимальных расстояний, а не просто наибольший путь (это же невозможно с циклами). То есть для данной $s$ нам нужно найти такую $t$, что $t-s$ максимально.

Беллман-Форд при запуске из $s$ найдет минимальные расстояния $s-x$, в том числе и $s-t$, но нам нужно $t-s$. Идея сделать инвертированный граф, то есть $\forall v-u$ ставим в инвертированый граф $u-v$, оставляя веса. Тогда у нас все так же будут отрицательные ребра без отрицательных циклов. Тогда заметим, что $w(s-t) = w(t-s)$. Тогда, запустив Беллмана Форда из $s$ в интвертированном графе, мы найдем минимальные расстояния от $s$ до всех $x$ инвертированного графа. Но мы знаем, что $d[s-x]$ - минимальное из всех путей в инвертированном графе; тогда $d[x-s]$ будет минимальным расстоянием из $x$ до $s$, так как вес пути в инвертированном графе равен весу пути в обычном графе. Тогда алгоритм поиска такой:

1) строим инвертированный граф за $O(E)$ - нужно развернуть каждое ребро
2) запускае Беллмана-Форда за $O(VE)$
3) среди получившихся $V$ расстояний $s-x$ ищем максимальное за $O(V)$ - нашли максимальное расстояние $s-t$ в инвертированном графе
4) это расстояние будет максимальным $t-s$ в нашем графе

Итого - $O(VE)$

\bigskip

\section*{\textbf{Задание 6}} \textbf{\textit{Кастелянша ходит по общежитию и производит устное внушение о последствиях ненадлежащего поведения живущих там студентов. Между двумя комнатами или можно пройти напрямую, или нельзя (в таком случае можно пройти через промежуточные комнаты). В общежитии $|V|$ комнат и $|E|$ коридоров между ними. Время прохождения каждого коридора известно и положительно. У одной комнаты может быть не более $k$ коридоров. Оказавшись в комнате, кастелянша всегда производит воспитательную работу, время которой составляет 3 минуты. Постройте алгоритм, позволяющий кастелянше по заданной конфигурации общежития быстро узнать, за какое минимальное время она может дойти от своего кабинета до заданной комнаты, докажите его корректность и оцените асимптотику.}}

Дан взвешенный неорграф G(V, E), веса путей неотрицательны. При этом на каждой вершине доп вес = 3, степень вершины не более k. 
Можно искать обычным алгоритмом Дейкстры, нам же нужен минимальный путь от вершины $s$ в пункт назначения $t$. Единственный вопрос - что делать с весами вершин. Просто в алгоритм релаксации мы делаем не так: $u.d = v.d + w(v, u)$, а  $u.d = v.d + w(v, u) + 3$, то есть как бы добавочныый вес ребру $v, u$ = времени, в течение которого мы будем находиться в комнате $v$. 

Корректность - обычный алгоритм Дейкстры, условие коридоров => нет отрицательных ребер (и соответственно циклов). Также заметим, что если путь $s - t$ состоит из $k$ ребер, то всего в комнатах мы потратим $3k$ минут: релаксация минимального пути включает $k$ релаксированных ребер, вследствие чего $u.d$ будет уже содержать эти $3k$ минут. Работает как обычный Дейкстра $O(|E| *lg|V|)$ (это обычно ребер больше чем вершин, но если полностью писать, то $O(|E| *lg(|V| + |E|))$ 

\bigskip


\section*{\textbf{Задание 7}} \textbf{\textit{Диаметр графа - это максимальное кратчайшее расстояниие между двумя его вершинами. Рассмотрим граф состояний кодового замка с 4 разрядами. Каждая комбинация соответствует одной вершине. Ребрами соединены комбинации, между которыми возможен переход вращением одного из колесиков на единицу вверх или вниз. Все ребра имеют единичный вес. Найдите диаметр этого графа.}}

Вершин будет $10^4$ - по одному на каждую комбинацию. Рассмотрим к-л вершину. У нее будет: $4 * 2$ ребер в соседние вершины. 

Рассмотрим комбинации $a_i, b_i, c_i, d_i$ и $a_j, b_j, c_j, d_j$, нам нужно покрутить замки из $i$ в $j$. Давайте рассмотрим пока что одно колесико.

Так как у нас циклично изменение, то максимальное количество поворотов будет равно 5, для противоположных частей цикла (0-5, 1-6 итд). Также заметим, что колесики у нас вращаются независимо, то есть такая логика корректна для каждого колесика: максимально необходимо 5 поворотов одного колесика, чтобы сменить состояние разряда на необходимое. То есть всего необходимо $5*4$ поворотов всех колесиков, чтобы сменить состояние кода. Так как один поворот = 1 ребру, то между любыми двумя вершинами существует путь, не превышающий 20. Отсюда диаметр графа = 20.

\bigskip

\end{document}